\newacronym{gnb}{GNB}{Grupo de Neurocomputación Biológica}
\newacronym{uam}{UAM}{Universidad Autónoma de Madrid}
\newacronym{gp}{GP}{\textit{Gaussian Process}}
\newacronym{se}{SE}{\textit{Squared Exponential}}
\newacronym{ard}{ARD}{\textit{Automatic Relevance Determination}}
\newacronym{ei}{EI}{\textit{Expected Improvement}}
\newacronym{lhs}{LHS}{\textit{Latin Hypercube Sampling}}
\newacronym{sos}{SOS}{\textit{Second-Order Sections}}
\newacronym{rt}{RT}{\textit{Real-Time}}
\newacronym{hrt}{HRT}{\textit{Hard \acl{rt}}}
\newacronym{nrt}{NRT}{\textit{No \acl{rt}}}
\newacronym{rtxi}{RTXI}{\textit{\acl{rt} eXperiment Interface}}
\newacronym{rthybrid}{RTHybrid}{\textit{\acl{rt} Hybrid}}
\newacronym{iir}{IIR}{\textit{Infinite Impulse Response}}
\newacronym{rprop}{Rprop}{\textit{Resilient Backpropagation}}
\newacronym{rk4}{RK4}{Runge-Kutta de cuarto orden}
\newacronym{rk5}{RK5}{Runge-Kutta de quinto orden}
\newacronym{daq}{DAQ}{\textit{Data Acquisition}}
\newacronym{gui}{GUI}{\textit{User Graphical Interface}}
\newacronym{pso}{PSO}{\textit{Particle Swarm Optimization}}
\newacronym{cmaes}{CMA-ES}{\textit{Covariance Matrix Adaptation Evolution Strategy}}
\newacronym{cpg}{CPG}{\textit{Central Pattern Generator}}
\newacronym{rtos}{RTOS}{\textit{\acl{rt} Operating System}}
\newacronym{bo}{BO}{\textit{Bayesian Optimization}}
\newacronym{hr}{HR}{Hindmarsh-Rose}


La construcción de circuitos neuronales computacionales y biohíbridos, en los que neuronas biológicas y modelos computacionales interactúan en tiempo real, \ac{rt}, exige sinapsis artificiales cuyos parámetros reproduzcan la dinámica de acoplamiento deseada. La parametrización de estos modelos sinápticos constituye un problema de optimización complejo, debido al elevado número de parámetros, su acoplamiento y la variabilidad inherente a cada preparación biológica, lo que hace la parametrización manual lenta e inasumible cuando el tiempo de vida de la preparación es limitado. El presente trabajo aborda este reto mediante un algoritmo de parametrización automática del modelo de sinapsis química de \cite{Golowasch1999}, adaptable a cada experimento y capaz de alcanzar dinámicas de acoplamiento específicas con comportamiento sináptico realista en un tiempo razonable.

\section{Motivación del trabajo\label{SEC:MOTIVACION}}

Los circuitos neuronales biohíbridos permiten estudiar la dinámica neuronal bajo condiciones controladas, conectando neuronas biológicas con modelos computacionales en \ac{rt} mediante sinapsis artificiales cuyos parámetros determinan la naturaleza, intensidad y temporalidad del acoplamiento.

Encontrar una parametrización adecuada para cada experimento no es trivial: cada preparación biológica presenta características electrofisiológicas distintas, por lo que los parámetros válidos en un experimento no son necesariamente trasladables a otro. Esta dificultad se acentúa con modelos detallados como el de sinapsis química de \cite{Golowasch1999}, cuyo espacio de parámetros presenta un notable acoplamiento que causa el fenómeno \textit{parameter sloppiness} \cite{Gutenkunst2007} (múltiples combinaciones producen respuestas similares). Dado que la parametrización debe repetirse en cada experimento, la lentitud del ajuste manual resulta perjudicial para preparaciones de vida limitada, por lo que se necesitan métodos de búsqueda automática eficientes y adaptados a esta problemática.

Este problema es de especial interés para las líneas de investigación del \ac{gnb} de la \ac{uam}. Motivados por esta necesidad experimental, el presente trabajo tiene como objetivo estudiar este modelo de sinapsis \cite{Golowasch1999} y desarrollar un algoritmo automático para parametrizarlo que se adapte a cada experimento para lograr dinámicas de acoplamiento específicas (fase o antifase) con un comportamiento sináptico realista y en un tiempo razonable.


\section{Formalización del problema\label{SEC:FORMALIZACION}}

Para contextualizar el problema que se aborda en este trabajo, esta sección presenta los conceptos teóricos fundamentales que sustentan la construcción de circuitos neuronales computacionales y biohíbridos y la optimización de sus parámetros.

\subsection{Modelo neuronal\label{SS:MODELONEURONAL}}

Un modelo neuronal es una descripción matemática que reproduce el comportamiento eléctrico de una neurona. Estos modelos capturan los patrones de actividad eléctrica característicos, como el disparo de potenciales de acción (\textit{spikes}) y las ráfagas de actividad (\textit{bursting}), en las que secuencias de \textit{spikes} se agrupan sobre una oscilación lenta del potencial de membrana. La elección del modelo neuronal implica un compromiso entre fidelidad biofísica y coste computacional \cite{Torres2012}: modelos detallados como el de Hodgkin-Huxley \cite{Hodgkin1952} describen con precisión las corrientes iónicas individuales, mientras que modelos fenomenológicos como el de \ac{hr} \cite{Hindmarsh1984} o el mapa de Rulkov \cite{Rulkov2002} utilizan formulaciones simplificadas que reproducen los patrones cualitativos de disparo con menor coste \cite{Fortuna2023}.

\subsection{Modelos de sinapsis\label{SS:MODELOSINAPSIS}}

Un modelo de sinapsis describe matemáticamente la transmisión de señales entre dos neuronas. Existen dos grandes categorías: los modelos de sinapsis eléctrica, que representan la conexión directa entre neuronas a través de uniones comunicantes \cite{GurelKazanci2007}, y los modelos de sinapsis química, que formalizan la transmisión mediada por neurotransmisores. Estos últimos ofrecen una mayor riqueza funcional, permitiendo capturar aspectos como el carácter excitatorio o inhibitorio de la conexión, la selectividad frecuencial \cite{TsodyksMarkram1998} o la dinámica temporal de activación de los canales sinápticos \cite{Destexhe1994}. En general, los parámetros del modelo (conductancias, potenciales de reversión, constantes cinéticas, entre otros) determinan conjuntamente las propiedades de la corriente generada y, por tanto, la naturaleza del acoplamiento entre las neuronas conectadas \cite{Golowasch1999}.

\subsection{Circuitos biohíbridos\label{SS:CIRCUITOBIOHIBRIDO}}

Un circuito biohíbrido conecta neuronas biológicas con modelos computacionales mediante sinapsis artificiales en \textit{software}. La técnica de \textit{Dynamic Clamp} \cite{Sharp1993a, Patel2017, Amaducci2019, Reyes-Sanchez2020} permite implementar esta interacción en lazo cerrado: el sistema adquiere el voltaje de la neurona viva, calcula los modelos sinápticos y neuronales, e inyecta la corriente resultante en la neurona viva de nuevo.

\subsection{\acl{hrt} (\acs{hrt})\label{SS:TIEMPOREAL}}

Las neuronas biológicas operan con regímenes temporales muy precisos, en escalas de milisegundos a segundos. Para que la interacción biohíbrida sea válida, el sistema de \textit{Dynamic Clamp} debe completar el cálculo de la corriente sináptica a inyectar en la neurona viva dentro de cada periodo de muestreo. Esto requiere ejecutar el software sobre un \ac{rtos} que garantice plazos temporales estrictos \cite{Buttazzo2011, Patel2017, Amaducci2019}. Esta restricción limita la complejidad computacional admisible de los modelos y algoritmos.

\subsection{Optimización de parámetros\label{SS:OPTIMIZACION}}

La parametrización automática de un modelo se formula como un problema de optimización: dado un conjunto de parámetros, se evalúa la calidad del resultado mediante una función objetivo que cuantifica la adecuación del comportamiento generado respecto al deseado. Un algoritmo de optimización explora el espacio paramétrico para encontrar la combinación que maximiza (o minimiza) esta función \cite{Nelder1965, Holland1975, Kennedy1995, Hansen2001, Shahriari2016, ReyesSnchez2023}.

\section{Estado del Arte\label{SEC:ESTADODELARTE}}

Esta sección revisa las principales herramientas y técnicas existentes relevantes para los circuitos biohíbridos: modelos de sinapsis, sistemas de \ac{rt} y algoritmos de optimización de caja negra.

\subsection{Modelos de sinapsis\label{SS:MODELOS_SINAPSIS}}

A continuación se presentan los principales modelos de sinapsis, ordenados de menor a mayor complejidad funcional.

\subsubsection{Sinapsis eléctricas}

La sinapsis eléctrica modela la conexión directa entre neuronas a través de uniones comunicantes (\textit{gap junctions}) \cite{GurelKazanci2007}, formalizándose como una conductancia lineal: $I_{\mathrm{gap}} = g_{\mathrm{gap}} \cdot (V_{\mathrm{post}} - V_{\mathrm{pre}})$. Ofrece simplicidad computacional y transmisión bidireccional instantánea, pero carece de la flexibilidad funcional de las sinapsis químicas: no distingue entre excitación e inhibición, no introduce filtrado frecuencial y no reproduce la riqueza dinámica (modulación, plasticidad, selectividad temporal) observada en la mayoría de conexiones químicas biológicas.

\subsubsection{Modelos cinéticos de sinapsis químicas}

Estos modelos describen la transmisión sináptica mediante ecuaciones que relacionan el potencial presináptico con la apertura de canales postsinápticos. El modelo de \cite{Destexhe1994} formula la cinética de unión del neurotransmisor al receptor como un sistema de primer orden. El de \cite{Golowasch1999}, empleado en este trabajo, separa la corriente en dos componentes (rápida y lenta) con cinéticas diferenciadas, capturando la transmisión mediada por \textit{spikes} (solos o junto con la onda lenta) y la transmisión gradual sobre la onda lenta. Ofrecen gran flexibilidad para reproducir comportamientos excitatorios e inhibitorios y ajustar independientemente la respuesta a distintas bandas frecuenciales, a costa de un mayor número de parámetros, algunos acoplados, que dificultan su ajuste manual.

\subsubsection{Modelos con plasticidad sináptica a corto plazo}

Los modelos de plasticidad a corto plazo, como el de Tsodyks-Markram \cite{TsodyksMarkram1998}, incorporan facilitación y depresión que modifican dinámicamente la eficacia sináptica según la actividad presináptica reciente, lo que permite capturar fenómenos como la adaptación a estímulos repetidos. Su complejidad añadida no resulta necesaria para este trabajo, cuyo objetivo es reproducir una dinámica de acoplamiento estacionaria sin requerir que la fuerza sináptica varíe con la historia de la actividad.

\subsubsection{Modelo de corrientes graduales rápida y lenta}

El modelo empleado en este trabajo es el de Golowasch et al. \cite{Golowasch1999}, diseñado para describir sinapsis graduales en el ganglio estomatogástrico de crustáceos. Su formulación con dos componentes (rápida y lenta) permite reconstruir modularmente la corriente sináptica aislando las contribuciones de alta frecuencia (\textit{spikes}) o la onda completa, y baja frecuencia  (onda lenta) de la señal presináptica. Esta separación frecuencial facilita la adaptación del acoplamiento a cada preparación experimental. Su formulación matemática se detalla en la Sección~\ref{SUBSEC:SINAPSIS_QUIMICA}.

\subsection{Sistemas y herramientas de \acl{rt}\label{SS:SISTEMAS_RT}}

La construcción de circuitos biohíbridos requiere plataformas con garantías de \ac{hrt}. \ac{rtxi} \cite{RTXI2010} es una plataforma de código abierto para \textit{Dynamic Clamp} que opera sobre un \ac{rtos} basado en un \textit{kernel} Linux parcheado con extensiones de \ac{hrt} como Xenomai \cite{Xenomai2004} (en este caso). Proporciona un entorno donde cada funcionalidad se implementa como un módulo independiente ejecutado periódicamente en el hilo de \ac{rt}, permitiendo componer circuitos combinando módulos. \ac{rthybrid} para \ac{rtxi} \cite{Amaducci2019} complementa esta plataforma con módulos para modelos neuronales, escalado de señales y análisis de \textit{bursts} \cite{Reyes-Sanchez2020}, y dispone de una versión \textit{standalone} independiente de \ac{rtxi}.

\subsection{Algoritmos de optimización de caja negra\label{SS:ALGORITMOS_OPTIMIZACION}}

Los algoritmos de optimización de caja negra buscan el óptimo de una función objetivo sin disponer de su expresión analítica ni de sus derivadas. A continuación se describen los principales métodos considerados para el problema, destacando sus ventajas e inconvenientes en este contexto.

\subsubsection{Método Nelder-Mead}

Nelder-Mead \cite{Nelder1965} es un algoritmo de búsqueda directa basado en transformaciones geométricas (reflexión, expansión, contracción y encogimiento) de un símplex en el espacio de parámetros. Ofrece simplicidad y bajo coste por iteración al no requerir información de gradientes, pero converge frecuentemente a óptimos locales, carece de garantías de convergencia global, su rendimiento se degrada mucho con alta dimensionalidad, y no dispone de mecanismos contra el \textit{parameter sloppiness}.

\subsubsection{Algoritmos genéticos}

Los algoritmos genéticos \cite{Holland1975} mantienen una población de soluciones candidatas y aplican selección, cruce y mutación para explorar el espacio de búsqueda. Son capaces de escapar de óptimos locales, pero requieren un número elevado de evaluaciones para converger, lo que resulta prohibitivo cuando cada evaluación es costosa en tiempo (como es el caso con neuronas vivas). No modelan la geometría del espacio de búsqueda ni las correlaciones entre parámetros, siendo ineficientes con \textit{parameter sloppiness} al no distinguir direcciones \textit{stiff} (donde pequeños cambios afectan significativamente al objetivo) y \textit{sloppy} (donde el objetivo es insensible) del espacio paramétrico, desperdiciando evaluaciones explorando combinaciones redundantes. En \cite{ReyesSnchez2023} se emplea uno para la parametrización de sinapsis biohíbridas, demostrando su viabilidad pero evidenciando la necesidad de diseñarlo cuidadosamente para obtener resultados aceptables con un presupuesto limitado de evaluaciones.

\subsubsection{\acl{pso} (\acs{pso})}

\ac{pso} \cite{Kennedy1995} simula el comportamiento colectivo de un enjambre donde cada partícula ajusta su posición según su mejor resultado individual y el global del grupo. Como método poblacional sin gradientes, ofrece buena exploración global pero comparte las limitaciones de los algoritmos genéticos: requiere muchas evaluaciones, es susceptible a convergencia prematura si la diversidad del enjambre se pierde y carece de mecanismos para modelar dependencias entre parámetros (malo con \textit{parameter sloppiness}).

\subsubsection{\acl{cmaes} (\acs{cmaes})}

\ac{cmaes} \cite{Hansen2001} es una estrategia evolutiva que adapta iterativamente la matriz de covarianza de una distribución normal multivariante para dirigir la búsqueda. Destaca por aprender la estructura de correlaciones: identifica direcciones \textit{stiff} y \textit{sloppy}, concentrando la búsqueda en las primeras, haciéndolo apto para el \textit{parameter sloppiness}. Sin embargo, requiere muchas evaluaciones para adaptar la matriz (coste creciente con la dimensionalidad), haciéndolo inadecuado para casos con neuronas vivas, donde cada evaluación es lenta.

\subsubsection{\acl{bo} (\acs{bo})}

La \ac{bo} \cite{Shahriari2016} construye un modelo probabilístico subrogado de la función objetivo (típicamente un \ac{gp} \cite{Rasmussen2006}) que se actualiza con cada evaluación. Una función de adquisición como \ac{ei} \cite{Mockus1978} selecciona el siguiente punto equilibrando exploración y explotación. Destaca en eficiencia muestral: al modelar globalmente la función objetivo y cuantificar la incertidumbre, encuentra buenas soluciones en pocas evaluaciones. Mediante \textit{kernels} con \ac{ard}, el \ac{gp} asigna escalas de longitud independientes a cada parámetro, identificando las direcciones relevantes y mitigando el \textit{parameter sloppiness}. Su coste computacional interno crece cúbicamente con las observaciones acumuladas, aunque este factor es secundario frente al coste de evaluación en experimentos con neuronas vivas. Como este es el método utilizado en este trabajo, se proporciona una explicación más detallada en la Sección~\ref{SUBSEC:OPTIMIZACION_BAYESIANA}.

\section{Objetivos\label{SEC:OBJETIVOS}}

A partir de los antecedentes expuestos, el objetivo principal de este trabajo es estudiar la parametrización del modelo de sinapsis química de \cite{Golowasch1999}, desarrollando un algoritmo que la automatice adaptándose a cada experimento para alcanzar dinámicas de acoplamiento específicas (fase o antifase) con un comportamiento sináptico realista. Los objetivos específicos son:

\begin{enumerate}
    \item Diseñar e implementar una función de evaluación que cuantifique la calidad de una parametrización sináptica considerando tanto la morfología como el rango de la corriente generada.
    \item Configurar y adaptar el algoritmo de optimización (\ac{bo}) al espacio de parámetros del modelo sináptico, incorporando mecanismos para mitigar el \textit{parameter sloppiness} inherente al problema.
    \item Desarrollar una implementación \textit{offline} que permita la parametrización rápida sobre señales pregrabadas usando una interacción unidireccional, facilitando la depuración y validación del algoritmo.
    \item Desarrollar un módulo \textit{online} integrado en \ac{rtxi} que ejecute la optimización en \ac{rt} durante el experimento, compatible con interacciones bidireccionales en modelos y en circuitos biohíbridos.
    \item Validar ambas implementaciones verificando que producen parametrizaciones que generan corrientes sinápticas con la morfología y escala esperadas en un tiempo razonable, así como comprobando la corrección del proceso de optimización.
\end{enumerate}
